\documentclass[11pt,a4paper]{article}
\usepackage[utf8]{inputenc}
\usepackage[margin=2.2cm]{geometry}
\usepackage{amsmath,amssymb,amsfonts,amsthm}
\usepackage{booktabs}
\usepackage{tabularx}
\usepackage[table]{xcolor}
\usepackage{listings}
\usepackage{fancyhdr}
\usepackage{hyperref}
\usepackage{tcolorbox}
\usepackage{enumitem}
\usepackage{microtype}

% Dynamically bind to the exact empirical LaTeX macros used by the IEEE manuscript
% Dynamic LaTeX Macros for Warehouse Digital Twin Optimization
% Auto-generated by generate_visualizations.py (Zero hardcoded values)

% Facility and Data Characteristics
\newcommand{\TotalLocations}{2,292}
\newcommand{\TotalOrderLines}{119,811}
\newcommand{\TotalOrders}{32,634}
\newcommand{\TotalPicks}{215,192}
\newcommand{\TotalWaves}{9,707}
\newcommand{\TotalOperationalDays}{178}
\newcommand{\BenchmarkWaves}{100}
\newcommand{\UnitGapReconciled}{268}
\newcommand{\AutomatedTestCount}{15}

% DaRA 2026 Timing Calibration
\newcommand{\DaRAValidGaps}{257}
\newcommand{\DaRAMedianGap}{36.0}
\newcommand{\DaRAMeanGap}{54.0}
\newcommand{\DaRAIQR}{26.0--59.0}
\newcommand{\CalibratedPickDwell}{12.6}
\newcommand{\DaRAHandleIQR}{9.1--20.6}
\newcommand{\CalibratedWalkSpeed}{1.10}

% Measured Per-Wave Optimizer Runtimes (ms)
\newcommand{\CASTWaveRuntimeMs}{2.56}
\newcommand{\CASTWaveRuntimeMaxMs}{6.34}
\newcommand{\StaticTSPRuntimeMs}{2.65}
\newcommand{\SShapeRuntimeMs}{0.04}

% Baseline FCFS + S-Shape
\newcommand{\MakespanSShapeMin}{7.26}
\newcommand{\DistSShapeM}{830.3}
\newcommand{\BlockedSShapeSec}{404.2}
\newcommand{\WorkloadCVSShape}{0.163}
\newcommand{\SShapeMakespanMean}{7.26}
\newcommand{\SShapeMakespanStd}{1.79}

% Baseline FCFS + Largest-Gap
\newcommand{\MakespanGapMin}{7.29}
\newcommand{\DistGapM}{831.1}
\newcommand{\BlockedGapSec}{406.7}
\newcommand{\WorkloadCVGap}{0.163}
\newcommand{\LargestGapMakespanMean}{7.29}
\newcommand{\LargestGapMakespanStd}{1.79}
\newcommand{\TStatVsLargestGap}{14.58}
\newcommand{\PValLargestGap}{< 10^{-25}}
\newcommand{\PValExactLargestGap}{2.16 \times 10^{-26}}
\newcommand{\WilcoxonWLargestGap}{27.0}
\newcommand{\WilcoxonPValLargestGap}{< 10^{-15}}
\newcommand{\CohenDLargestGap}{-1.46}

% Baseline FCFS + Static TSP
\newcommand{\MakespanTSPMin}{6.93}
\newcommand{\DistTSPM}{807.7}
\newcommand{\BlockedTSPSec}{380.0}
\newcommand{\WorkloadCVTSP}{0.147}
\newcommand{\StaticTSPMakespanMean}{6.93}
\newcommand{\StaticTSPMakespanStd}{1.70}
\newcommand{\TStatVsStaticTSP}{13.42}
\newcommand{\PValStaticTSP}{< 10^{-23}}
\newcommand{\PValExactStaticTSP}{5.25 \times 10^{-24}}
\newcommand{\WilcoxonWStaticTSP}{67.0}
\newcommand{\WilcoxonPValStaticTSP}{< 10^{-14}}
\newcommand{\CohenDVsStaticTSP}{-1.34}

% Proposed CAST-Wave
\newcommand{\MakespanCASTMin}{5.32}
\newcommand{\DistCASTM}{624.8}
\newcommand{\BlockedCASTSec}{158.6}
\newcommand{\WorkloadCVCAST}{0.212}
\newcommand{\CASTWaveMakespanMean}{5.32}
\newcommand{\CASTWaveMakespanStd}{1.18}

% Sector + Static TSP (Variant 4)
\newcommand{\MakespanSectorMin}{5.37}
\newcommand{\MakespanSectorStd}{1.21}
\newcommand{\SectorReductionVsFairPct}{26.1}

% Relative Improvements
\newcommand{\MakespanRedSShapePct}{26.7}
\newcommand{\MakespanRedLargestGapPct}{27.0}
\newcommand{\MakespanRedTSPPct}{23.2}
\newcommand{\MakespanImprovementVsSShape}{26.7}
\newcommand{\MakespanImprovementVsLargestGap}{27.0}
\newcommand{\MakespanImprovementVsStaticTSP}{23.2}
\newcommand{\DistanceImprovementVsSShape}{24.8}
\newcommand{\DistanceImprovementVsStaticTSP}{22.7}
\newcommand{\BlockedTimeCASTWaveMean}{158.6}
\newcommand{\BlockedTimeFCFSMean}{380.0}
\newcommand{\BlockedTimeReductionPct}{60.8}
\newcommand{\BlockedTimeReductionVsTSPPct}{58.3}

% Workload CV Context
\newcommand{\WorkloadCVLegacy}{1.732}
\newcommand{\WorkloadCVFairRR}{0.163}
\newcommand{\WorkloadCVLPT}{0.147}
\newcommand{\WorkloadCVImprovementVsLegacyPct}{87.8}

% Statistical Tests vs Fair S-Shape
\newcommand{\TStatSShape}{-14.26}
\newcommand{\PValSShape}{< 10^{-25}}
\newcommand{\PValExactSShape}{9.73 \times 10^{-26}}
\newcommand{\CohensDSShape}{-1.43}
\newcommand{\TStatTSP}{-13.42}
\newcommand{\PValTSP}{< 10^{-23}}
\newcommand{\CohensDTSP}{-1.34}
\newcommand{\CohenDVsSShape}{-1.43}
\newcommand{\TStatVsSShape}{14.26}
\newcommand{\WilcoxonWSShape}{36.0}
\newcommand{\WilcoxonPValSShape}{< 10^{-15}}
\newcommand{\WilcoxonPValExactSShape}{1.78 \times 10^{-16}}
\newcommand{\WinsVsSShape}{90}
\newcommand{\LossesVsSShape}{3}
\newcommand{\TiesVsSShape}{7}

% Sector vs CAST Comparison
\newcommand{\WinsVsSector}{19}
\newcommand{\LossesVsSector}{12}
\newcommand{\TiesVsSector}{69}
\newcommand{\TStatSectorVsCAST}{1.10}
\newcommand{\PValSectorVsCAST}{0.273}
\newcommand{\WilcoxonWSectorVsCAST}{188.0}
\newcommand{\WilcoxonPValSectorVsCAST}{0.240}

% Dense Waves
\newcommand{\DenseWavesCount}{20}
\newcommand{\DenseMeanPicks}{46.6}
\newcommand{\DenseSectorBetterFleets}{4, 6, 8, 10, 12}
\newcommand{\DenseCASTBetterFleets}{2}
\newcommand{\DenseMkFCFSTwo}{15.25}
\newcommand{\DenseMkSectorTwo}{13.38}
\newcommand{\DenseMkCASTTwo}{13.32}
\newcommand{\DenseBlockedFCFSTwo}{274.7}
\newcommand{\DenseBlockedSectorTwo}{64.9}
\newcommand{\DenseBlockedCASTTwo}{51.7}
\newcommand{\DenseRedCASTvsFCFSTwo}{12.6}
\newcommand{\DenseRedSectorvsFCFSTwo}{12.3}
\newcommand{\DenseMkFCFSFour}{11.26}
\newcommand{\DenseMkSectorFour}{8.47}
\newcommand{\DenseMkCASTFour}{8.48}
\newcommand{\DenseBlockedFCFSFour}{593.2}
\newcommand{\DenseBlockedSectorFour}{160.8}
\newcommand{\DenseBlockedCASTFour}{154.2}
\newcommand{\DenseRedCASTvsFCFSFour}{24.7}
\newcommand{\DenseRedSectorvsFCFSFour}{24.8}
\newcommand{\DenseMkFCFSSix}{9.71}
\newcommand{\DenseMkSectorSix}{6.90}
\newcommand{\DenseMkCASTSix}{7.09}
\newcommand{\DenseBlockedFCFSSix}{835.3}
\newcommand{\DenseBlockedSectorSix}{298.0}
\newcommand{\DenseBlockedCASTSix}{279.9}
\newcommand{\DenseRedCASTvsFCFSSix}{27.0}
\newcommand{\DenseRedSectorvsFCFSSix}{28.9}
\newcommand{\DenseMkFCFSEight}{8.50}
\newcommand{\DenseMkSectorEight}{6.26}
\newcommand{\DenseMkCASTEight}{6.37}
\newcommand{\DenseBlockedFCFSEight}{1000.3}
\newcommand{\DenseBlockedSectorEight}{416.3}
\newcommand{\DenseBlockedCASTEight}{402.5}
\newcommand{\DenseRedCASTvsFCFSEight}{25.1}
\newcommand{\DenseRedSectorvsFCFSEight}{26.4}
\newcommand{\DenseMkFCFSTen}{7.85}
\newcommand{\DenseMkSectorTen}{5.42}
\newcommand{\DenseMkCASTTen}{5.67}
\newcommand{\DenseBlockedFCFSTen}{1165.2}
\newcommand{\DenseBlockedSectorTen}{463.6}
\newcommand{\DenseBlockedCASTTen}{476.2}
\newcommand{\DenseRedCASTvsFCFSTen}{27.7}
\newcommand{\DenseRedSectorvsFCFSTen}{31.0}
\newcommand{\DenseMkFCFSTwelve}{7.01}
\newcommand{\DenseMkSectorTwelve}{5.13}
\newcommand{\DenseMkCASTTwelve}{5.43}
\newcommand{\DenseBlockedFCFSTwelve}{1273.3}
\newcommand{\DenseBlockedSectorTwelve}{548.7}
\newcommand{\DenseBlockedCASTTwelve}{593.7}
\newcommand{\DenseRedCASTvsFCFSTwelve}{22.6}
\newcommand{\DenseRedSectorvsFCFSTwelve}{26.8}

% Component ablation (every metric measured per stage)
\newcommand{\LegacyMakespanMin}{11.74}
\newcommand{\LegacyMakespanStd}{2.39}
\newcommand{\LegacyIncreaseVsFairPct}{61.6}
\newcommand{\StaticTSPRedVsSShapePct}{4.6}
\newcommand{\LPTIdenticalWaves}{100}
\newcommand{\WorkloadCVFairTSP}{0.147}
\newcommand{\WorkloadCVSector}{0.211}
\newcommand{\DistSectorM}{621.4}
\newcommand{\BlockedSectorSec}{168.7}
\newcommand{\SectorShareOfGainPct}{98}


% Adjust header height for fancyhdr
\setlength{\headheight}{14.5pt}
\addtolength{\topmargin}{-2.5pt}

% Define color palette
\definecolor{primary}{RGB}{18, 52, 86}     % Deep Navy
\definecolor{secondary}{RGB}{204, 85, 0}   % Burnt Orange / Amber
\definecolor{darkgray}{RGB}{50, 50, 50}
\definecolor{lightbg}{RGB}{248, 249, 250}
\definecolor{codebg}{RGB}{244, 245, 247}
\definecolor{codegreen}{RGB}{34, 139, 34}
\definecolor{codeblue}{RGB}{25, 25, 112}
\definecolor{codepurple}{RGB}{128, 0, 128}

\hypersetup{
    colorlinks=true,
    linkcolor=primary,
    citecolor=primary,
    urlcolor=secondary
}

% Page styling
\pagestyle{fancy}
\fancyhf{}
\rhead{\textcolor{darkgray}{\footnotesize Industrial Warehouse Digital Twin: Theory, Code Architecture \& Model Rationale}}
\lhead{\textcolor{darkgray}{\footnotesize Technical Reference Document}}
\rfoot{\thepage}
\renewcommand{\headrulewidth}{0.4pt}

% Code listing styling
\lstdefinestyle{pythonstyle}{
    backgroundcolor=\color{codebg},
    commentstyle=\color{codegreen}\itshape,
    keywordstyle=\color{codeblue}\bfseries,
    numberstyle=\tiny\color{darkgray},
    stringstyle=\color{codepurple},
    basicstyle=\ttfamily\footnotesize,
    breakatwhitespace=false,
    breaklines=true,
    captionpos=b,
    keepspaces=true,
    numbers=left,
    numbersep=6pt,
    showspaces=false,
    showstringspaces=false,
    showtabs=false,
    tabsize=4,
    frame=single,
    rulecolor=\color{lightgray}
}
\lstset{style=pythonstyle}

% Custom boxes
\newtcolorbox{callout}[2][]{
    colback=lightbg,
    colframe=primary,
    fonttitle=\bfseries\color{white},
    colbacktitle=primary,
    title=#2,
    arc=2mm,
    #1
}

\newtcolorbox{rationalebox}[2][]{
    colback=lightbg,
    colframe=secondary,
    fonttitle=\bfseries\color{white},
    colbacktitle=secondary,
    title=#2,
    arc=2mm,
    #1
}

\newtheorem{proposition}{Proposition}

\title{\textbf{\LARGE Theoretical Foundations, Block-by-Block Codebase Architecture,\\Data Cleaning Pipeline, and Model Selection Rationale}\\[0.5em]
\large An Empirical Warehouse Digital Twin and Spatiotemporal Wave Optimizer}
\author{\textbf{Cristhian David Caceres Mateus}, \textbf{Sharique Haroon}, \textbf{Usaid Aamer}, \textbf{Grejs Ibrahimi}, and \textbf{Shan Faiz}\textsuperscript{*}\\[0.4em]
\normalsize Department of Business, University of Europe for Applied Sciences, Potsdam 14469, Germany\\
\normalsize Email: \texttt{\{cristhian.caceres, sharique.haroon, usaid.aamer, grejs.ibrahimi, shan.faiz\}@ue-germany.de}\\
\normalsize\textsuperscript{*}Faculty Supervisor}
\date{November 2026}

\begin{document}

\maketitle

\begin{abstract}
This technical reference document provides an exhaustive theoretical and block-by-block engineering breakdown of the complete data cleaning, discrete-event simulation, optimization, and statistical evaluation pipeline in \texttt{WarehouseDatasets2}. Every numerical metric in this document is dynamically bound via \texttt{paper/tables/macros.tex} to the verified empirical CSV artifacts. First, we present the mathematical theory governing the 44-node topological navigation graph, the continuous-time sub-aisle congestion physics, the three-stage \textbf{CAST-Wave} (Congestion-Aware Spatio-Temporal Batching and Wave Scheduler) algorithm, and the parametric and non-parametric statistical hypothesis tests (paired Student's $t$-test, Wilcoxon signed-rank test, Cohen's $d_z$ effect size, and win/loss/tie dominance analysis). Second, we walk block-by-block through the two-stage data hygiene workflow (\texttt{src/clean\_data.py} across four 2025--2026 intralogistics datasets and \texttt{src/data\_preprocessing.py} for coordinate reconciliation), the simulation and optimization modules, and the \AutomatedTestCount-test automated verification suite. Finally, we provide a rigorous comparative rationale for selecting CAST-Wave over Mixed-Integer Linear Programming (MILP), Multi-Agent Reinforcement Learning (MARL), Genetic Algorithms (GA), and unclustered TSP heuristics.
\end{abstract}

\vspace{1em}
\tableofcontents
\newpage

% ============================================================================
% SECTION 1: ARCHITECTURE OVERVIEW
% ============================================================================
\section{System Overview and End-to-End Architecture}

The digital twin and optimization framework in \texttt{WarehouseDatasets2} evaluates internal logistics workflows within a footwear manufacturing distribution facility comprising \TotalLocations~active storage locations across 18 narrow aisles ($A$--$R$), two storage blocks (West and East), and three vertical cross-corridors (\texttt{LC}, \texttt{CC}, \texttt{RC}). The software architecture integrates the team's multi-dataset raw auditing pipeline with continuous-time simulation physics, wave optimization, and LaTeX-synchronized evaluation:

\begin{figure}[h!]
\centering
\begin{callout}{End-to-End Pipeline Data Flow}
\begin{enumerate}[leftmargin=*]
    \item \textbf{Stage 1 --- Multi-Dataset Raw Auditing \& Cleaning} (\texttt{src/clean\_data.py} \& \texttt{clean\_data.py}): Cleans all four raw 2025--2026 datasets in \texttt{data/01\_raw/} (Footwear Warehouse WMS, KIT RADAR4KIT Benchmark, Fraunhofer IML DaRA wearable timing, and RMFS robotic pods); strips UTF-8 BOMs and mixed delimiters; drops 2,559 exact 9-field duplicate order lines; classifies shoe size scales; builds the 44-node, 68-edge navigation graph; and outputs 23 clean CSV/JSON artifacts into \texttt{data/clean/} plus \texttt{CLEANING\_LOG.md}.
    \item \textbf{Stage 2 --- Coordinate Reconciliation \& Parquet Indexing} (\texttt{src/data\_preprocessing.py}): Resolves the open modelling checklist from \texttt{CLEANING\_LOG.md}; maps the 22 unlisted staging/overflow and extended-bay codes (\texttt{RC-01}..\texttt{RC-05}, \texttt{EX-01}..\texttt{EX-11}, \texttt{C-014}..\texttt{C-102}, \texttt{ZN-COR}, and \texttt{E-25-20}..\texttt{K-26-20}) to metric coordinates; links SKU ABC classes; and exports validated Parquet tables to \texttt{data/02\_processed/}.
    \item \textbf{Stage 3 --- Topological Graph Engine} (\texttt{src/digital\_twin/warehouse\_graph.py}): Loads the 44-node, 68-edge metric network; establishes the central I/O depot at node \texttt{CC-08} $(32.24\text{ m}, 50.48\text{ m})$; and precomputes all-pairs Dijkstra shortest-path distance matrices.
    \item \textbf{Stage 4 --- Discrete-Event Simulator} (\texttt{src/digital\_twin/simulation\_engine.py}): Continuous-time event engine tracking kinematic motion ($v_0 = \CalibratedWalkSpeed\text{ m/s}$), narrow sub-aisle conflict physics ($K_{\text{aisle}}=1$ picker limit per West/East aisle segment, $\alpha=0.40$ speed penalty), and empirical DaRA pick handling dwell ($\tau_{\text{pick}} = \CalibratedPickDwell\text{ s}$).
    \item \textbf{Stage 5 --- Wave Optimization} (\texttt{src/optimization/}): 
    \begin{itemize}
        \item Baselines: Fair Multi-Picker FCFS (Round-Robin) + S-Shape, Fair FCFS + Largest-Gap, Fair FCFS + Static TSP (2-Opt), Balanced LPT + Static TSP, and Legacy Single-Picker FCFS ($P_{\text{eff}}=1$).
        \item Proposed: \textbf{CAST-Wave} (5-zone topological sector clustering, capacity-constrained batch allocation with $C = 27\text{ units}$, and dynamic congestion-penalized 2-Opt tour sequencing with $\lambda = 15.0\text{ m}$).
    \end{itemize}
    \item \textbf{Stage 6 --- Evaluation \& Macro Automation} (\texttt{src/evaluation/}): Runs \BenchmarkWaves~historical wave benchmarks, 5-stage component ablations, and 3-arm dense-wave fleet scalability ablations ($P \in \{2, 4, 6, 8, 10, 12\}$ across \DenseWavesCount~merged waves averaging $\DenseMeanPicks$~picks); writes dynamic metrics to \texttt{paper/tables/macros.tex} and generates vector graphics.
\end{enumerate}
\end{callout}
\end{figure}

% ============================================================================
% SECTION 2: THEORETICAL FOUNDATIONS & STATISTICAL METHODOLOGY
% ============================================================================
\newpage
\section{Mathematical Theory of the Model and Statistical Validation}
\label{sec:theory}

\subsection{Graph-Theoretic Warehouse Manifold and Shortest-Path Metric}
In a physical multi-block warehouse, storage racks form impassable barriers that invalidate unconstrained Euclidean ($L_2$) or rectilinear ($L_1$) distance metrics. Specifically, aisles $A$--$G$ and $R$ terminate at the central cross-corridor (\texttt{CC}, $x = 32.24\text{ m}$), whereas aisles $H$--$Q$ extend across both the West block ($x \in [5.28, 32.24]\text{ m}$) and the East block ($x \in [32.24, 54.88]\text{ m}$).

We model the navigable floorplan as an undirected weighted graph $\mathcal{G} = (\mathcal{V}, \mathcal{E}, w)$, where $|\mathcal{V}| = 44$ vertices represent corridor-aisle intersections along the Left Corridor (\texttt{LC}, $x = 5.28\text{ m}$), Center Corridor (\texttt{CC}, $x = 32.24\text{ m}$, housing the I/O depot \texttt{CC-08} at $y = 50.48\text{ m}$), and Right Corridor (\texttt{RC}, $x = 54.88\text{ m}$), and $|\mathcal{E}| = 68$ edges represent walkable corridor and aisle segments with metric length $w(e) = \| \mathbf{x}_u - \mathbf{x}_v \|_2$.

Each pick task $\tau_i$ at storage coordinate $\mathbf{p}_i = (x_i, y_i, a_i)$ resides in aisle $a_i$ with centerline ordinate $y_c(a_i)$ and belongs to horizontal block:
\begin{equation}
    b(x_i) = \begin{cases}
        \text{west}, & \text{if } x_i \le 32.24\text{ m}, \\
        \text{east}, & \text{if } x_i > 32.24\text{ m}.
    \end{cases}
\end{equation}
The exact walking distance $d(\mathbf{p}_1, \mathbf{p}_2)$ between two pick locations is governed by two cases:
\begin{enumerate}[leftmargin=*]
    \item \textbf{Intra-Segment Traversal} ($a_1 = a_2$ and $b(x_1) = b(x_2)$): The operator walks directly along the aisle centerline without entering a vertical cross-corridor:
    \begin{equation}
        d(\mathbf{p}_1, \mathbf{p}_2) = |x_1 - x_2| + |y_1 - y_c(a_1)| + |y_2 - y_c(a_2)|.
    \end{equation}
    \item \textbf{Inter-Segment Traversal} ($a_1 \neq a_2$ or $b(x_1) \neq b(x_2)$): Let $\mathcal{X}(\mathbf{p}_1) \subset \mathcal{V}$ and $\mathcal{E}(\mathbf{p}_2) \subset \mathcal{V}$ denote the two bounding corridor intersection nodes of segment $(a_1, b(x_1))$ and segment $(a_2, b(x_2))$, respectively, and let $d_{\mathcal{G}}(u, v)$ denote the precomputed all-pairs Dijkstra shortest path in $\mathcal{G}$. The exact transition distance is the minimum over the four corridor exit/entry combinations:
    \begin{equation}
        d(\mathbf{p}_1, \mathbf{p}_2) = \min_{u \in \mathcal{X}(\mathbf{p}_1), \, v \in \mathcal{E}(\mathbf{p}_2)} \Big[ \big(|x_1 - x_u| + |y_1 - y_u|\big) + d_{\mathcal{G}}(u, v) + \big(|x_2 - x_v| + |y_2 - y_v|\big) \Big].
    \end{equation}
\end{enumerate}

\subsection{Continuous-Time Discrete-Event Kinematics and Congestion Physics}
Let $\mathcal{P} = \{0, 1, \dots, P-1\}$ denote the active picker fleet assigned to a wave with ordered tour sequences $\{\pi_0, \dots, \pi_{P-1}\}$. The digital twin advances continuous simulation time $t \ge 0$ via a min-heap priority queue $\mathcal{Q}$ of events $e = (t, p, \text{type}, \text{step})$.

\subsubsection{Sub-Aisle Contention and Blocking Delay Derivation}
Each half-aisle sub-segment $s = (a, b) \in \{A, \dots, R\} \times \{\text{west}, \text{east}\}$ maintains an active picker occupancy set $\Omega(s, t) \subseteq \mathcal{P}$. Because narrow footwear aisles restrict unimpeded passing to $K_{\text{aisle}} = 1$ trolley, a picker $p$ entering segment $s$ at time $t$ experiences unimpeded velocity $v_0 = \CalibratedWalkSpeed\text{ m/s}$ if $|\Omega(s, t) \setminus \{p\}| < K_{\text{aisle}}$, and throttled velocity $v_{\text{cong}} = \alpha v_0$ ($\alpha = 0.40$) whenever $|\Omega(s, t) \setminus \{p\}| \ge K_{\text{aisle}}$.

For a travel leg of distance $d$ entering an occupied sub-segment, the unimpeded travel duration is $t_{\text{free}} = d / v_0$, whereas the congested traversal duration is $t_{\text{cong}} = d / (\alpha v_0)$. The induced physical blocking delay $\Delta t_{\text{blk}}$ is therefore:
\begin{equation}
    \Delta t_{\text{blk}}(d) = t_{\text{cong}} - t_{\text{free}} = \frac{d}{\alpha v_0} - \frac{d}{v_0} = \frac{d}{v_0}\left(\frac{1 - \alpha}{\alpha}\right).
\end{equation}
With $\alpha = 0.40$, $(1 - \alpha)/\alpha = 1.50$, meaning sub-aisle contention inflates traversal time on that leg by $150\%$ ($\approx 12\text{--}15\text{ s}$ on a typical $9\text{--}11\text{ m}$ aisle entry leg), plus the empirical handling dwell $\tau_{\text{pick}} = \CalibratedPickDwell\text{ s}$.

\subsubsection{Wave-Level Performance Metrics and Analytical Proofs}
Once all $P$ pickers complete their tours and return to the depot \texttt{CC-08}, let $T_p = t_{\text{travel}, p} + t_{\text{pick}, p}$ denote the total active tour duration of picker $p$. The four primary performance metrics are:
\begin{align}
    \text{Makespan:} \quad C_{\max} &= \max_{p \in \mathcal{P}} T_p, \\
    \text{Fleet Travel Distance:} \quad D_{\text{fleet}} &= \sum_{p \in \mathcal{P}} d_p, \\
    \text{Cumulative Blocking Time:} \quad T_{\text{blk}} &= \sum_{p \in \mathcal{P}} t_{\text{blocked}, p}, \\
    \text{Workload CV:} \quad \text{CV}_W &= \frac{\sigma_T}{\mu_T} = \frac{\sqrt{\frac{1}{P}\sum_{p=0}^{P-1} (T_p - \bar{T})^2}}{\frac{1}{P}\sum_{p=0}^{P-1} T_p}.
\end{align}

Two key structural properties observed in the empirical experiments follow directly from these definitions:

\begin{proposition}[Single-Picker Starvation CV Bound]
Under legacy sequential FCFS batching without round-robin coordination, whenever total wave demand satisfies $\sum_{i=1}^M q_i \le C = 27$, Picker $0$ receives all $M$ tasks ($T_0 > 0$) while Pickers $1, \dots, P-1$ receive zero tasks ($T_1 = \dots = T_{P-1} = 0$). The resulting Workload Coefficient of Variation is identically $\text{CV}_W = \sqrt{P - 1}$, independent of $T_0$.
\end{proposition}
\begin{proof}
With $T_0 > 0$ and $T_1 = \dots = T_{P-1} = 0$, the fleet mean is $\bar{T} = T_0 / P$. The population variance across the $P$ pickers is:
\begin{equation}
    \sigma_T^2 = \frac{1}{P}\left[ \left(T_0 - \frac{T_0}{P}\right)^2 + (P - 1)\left(0 - \frac{T_0}{P}\right)^2 \right] = \frac{T_0^2}{P}\left[ \frac{(P-1)^2}{P^2} + \frac{P-1}{P^2} \right] = \frac{T_0^2 (P - 1)}{P^2}.
\end{equation}
Taking the square root and dividing by $\bar{T} = T_0 / P$ yields $\text{CV}_W = \sigma_T / \bar{T} = \sqrt{P - 1}$. For $P = 4$, $\text{CV}_W = \sqrt{3} \approx 1.732051$, matching the empirical legacy column in Table~\ref{tab:doc_component_ablation} to six decimal places.
\end{proof}

\begin{proposition}[Equivalence of Unit LPT and Round-Robin Under Single-Pair Picks]
When $q_i = 1$ for all $i \in \{1, \dots, M\}$, discrete Longest Processing Time (LPT) unit bin packing produces the exact same batch partition as cyclic Round-Robin assignment.
\end{proposition}
\begin{proof}
In \texttt{picks\_clean.parquet}, $100\%$ of the \TotalPicks~pick records specify $q_i = 1$ pair of shoes. Stable sorting of tasks in non-increasing order of $q_i$ preserves their original index sequence $(0, 1, \dots, M-1)$. At step $i$, picker loads satisfy $L_p = \lfloor i / P \rfloor + \mathbb{I}(p < i \bmod P)$, so the unique least-loaded picker with lowest index is $p^* = i \bmod P$, which is identical to cyclic Round-Robin dispatch.
\end{proof}

\subsection{Theoretical Mechanics of the CAST-Wave Three-Stage Optimizer}
CAST-Wave overcomes the combinatorial intractability of joint batching and time-indexed routing by decomposing wave scheduling into three polynomial-time stages:

\begin{enumerate}[leftmargin=*]
    \item \textbf{Stage 1 --- Topological 5-Zone Sector Decomposition}:
    Each pick task $\tau_i = (x_i, y_i)$ is mapped to one of five topological zones aligned with the Center Corridor ($x = 32.24\text{ m}$) and longitudinal aisle groups:
    \begin{equation}
        \sigma(x_i, y_i) = \begin{cases}
            0 \quad (\text{West South: Aisles } A\text{--}D), & x_i \le 32.24, \; y_i \le 21.6, \\
            1 \quad (\text{West Mid: Aisles } E\text{--}G), & x_i \le 32.24, \; 21.6 < y_i \le 42.4, \\
            2 \quad (\text{West North: Aisles } H\text{--}R), & x_i \le 32.24, \; y_i > 42.4, \\
            3 \quad (\text{East Mid: Aisles } H\text{--}L), & x_i > 32.24, \; y_i \le 76.0, \\
            4 \quad (\text{East North: Aisles } M\text{--}Q), & x_i > 32.24, \; y_i > 76.0.
        \end{cases}
    \end{equation}

    \item \textbf{Stage 2 --- Sector-Cohesive Capacity-Constrained Batching}:
    Within each sector $s \in \{0, \dots, 4\}$, tasks are ordered by depot distance $d(\text{CC-08}, \tau_i)$. Let $L_{\text{target}} = \lceil \sum_{i=1}^M q_i / P \rceil$ denote the balanced unit target and $\mathcal{S}_p \subseteq \{0, \dots, 4\}$ the set of sectors already assigned to picker $p$. For each task $\tau_i$ in sector $s$, among pickers with $L_p < L_{\text{target}}$ and $L_p + q_i \le C = 27$, CAST-Wave selects:
    \begin{equation}
        p^* = \arg\min_{p} \Big( \mathbb{I}[s \notin \mathcal{S}_p], \; L_p \Big).
    \end{equation}
    Prioritizing $\mathbb{I}[s \notin \mathcal{S}_p] = 0$ keeps all tasks of sector $s$ consolidated onto the same operator until their target load $L_{\text{target}}$ is reached, preventing multiple pickers from being sent into the same topological zone.

    \item \textbf{Stage 3 --- Spatiotemporal Reservation Table and Congestion-Aware 2-Opt}:
    As each picker $p \in \{0, \dots, P-1\}$ is routed, their predicted sub-aisle occupancy intervals $[t_{\text{arr}}^{(k)}, t_{\text{dep}}^{(k)}]$ (where $t_{\text{dep}}^{(k)} = t_{\text{arr}}^{(k)} + \tau_{\text{pick}}$) are logged in a shared reservation table $\mathcal{O}(a, b)$. For subsequent pickers, the overlap count at candidate arrival time $t_{\text{arr}}$ in sub-segment $(a, b)$ is:
    \begin{equation}
        \Phi(a, b, t_{\text{arr}}) = \sum_{[t_s, t_e] \in \mathcal{O}(a, b)} \mathbb{I}\Big( \max(t_{\text{arr}}, t_s) \le \min(t_{\text{arr}} + \tau_{\text{pick}}, t_e) \Big).
    \end{equation}
    Both the initial Nearest-Neighbor tour construction and the subsequent 2-Opt local search refinement evaluate candidate tours $\pi_p = (\tau_{(1)}, \dots, \tau_{(m_p)})$ using the penalized total tour objective:
    \begin{equation}
        J_{\text{total}}(\pi_p) = d(\text{CC-08}, \tau_{(1)}) + \sum_{k=1}^{m_p - 1} d(\tau_{(k)}, \tau_{(k+1)}) + d(\tau_{(m_p)}, \text{CC-08}) + \lambda \sum_{k=1}^{m_p} \Phi\big(a_{(k)}, b_{(k)}, t_{\text{arr}}^{(k)}(\pi_p)\big),
    \end{equation}
    where $\lambda = 15.0\text{ m}$ equals the unimpeded walking distance corresponding to a $\approx 13.6\text{ s}$ sub-aisle yielding delay ($\lambda = v_0 \cdot \Delta t_{\text{blk}}$). In each 2-Opt iteration, reversing sub-segment $\pi_p[i:j]$ is accepted if and only if $J_{\text{total}}(\pi_p^{\text{new}}) < J_{\text{total}}(\pi_p) - 10^{-4}$, ensuring local search eliminates geometric crossovers without undoing valid congestion detours.
\end{enumerate}

\subsection{Mathematical Theory of the Statistical Hypothesis Tests}
\label{sec:statistical_theory}
To rigorously validate whether CAST-Wave's makespan reductions are statistically significant and practically meaningful across heterogeneous demand profiles, we employ a four-pillar statistical framework across all $N = \BenchmarkWaves$ historical waves.

\subsubsection{1. Paired Randomized Block Design}
Because historical waves vary substantially in order line count ($16$ to $30$ picks) and spatial dispersion across ABC classes, treating two policies as independent samples would confound algorithmic performance with inter-wave variance $\sigma_{\text{wave}}^2$. By evaluating both policies on the exact same wave $i \in \{1, \dots, N\}$, we analyze the paired intra-wave difference:
\begin{equation}
    d_i = T_i^{\text{CAST}} - T_i^{\text{Base}}, \quad i = 1, \dots, N.
\end{equation}
By the variance of linear combinations, pairing removes the positive covariance $\operatorname{Cov}(T^{\text{CAST}}, T^{\text{Base}}) > 0$ induced by wave size:
\begin{equation}
    \sigma_d^2 = \operatorname{Var}(T^{\text{CAST}} - T^{\text{Base}}) = \sigma_{\text{CAST}}^2 + \sigma_{\text{Base}}^2 - 2\operatorname{Cov}(T^{\text{CAST}}, T^{\text{Base}}) \ll \sigma_{\text{CAST}}^2 + \sigma_{\text{Base}}^2,
\end{equation}
providing maximum statistical power to isolate the true algorithmic effect $\mu_d = \mathbb{E}[d_i]$.

\subsubsection{2. Parametric Paired Student's $t$-Test}
\begin{itemize}[leftmargin=*]
    \item \textbf{Mathematical Formulation}: Let $\bar{d}$ and $s_d$ denote the sample mean and sample standard deviation of the $N$ paired differences:
    \begin{equation}
        \bar{d} = \frac{1}{N}\sum_{i=1}^N d_i, \qquad s_d = \sqrt{\frac{1}{N - 1}\sum_{i=1}^N (d_i - \bar{d})^2}, \qquad \text{SE}_{\bar{d}} = \frac{s_d}{\sqrt{N}}.
    \end{equation}
    To test the null hypothesis $H_0: \mu_d = 0$ against the two-sided alternative $H_1: \mu_d \neq 0$, the paired $t$-statistic is:
    \begin{equation}
        t = \frac{\bar{d}}{\text{SE}_{\bar{d}}} = \frac{\bar{d}}{s_d / \sqrt{N}} \sim t_{(\nu = N - 1)},
    \end{equation}
    with $\nu = 99$ degrees of freedom and two-sided $p$-value $p = 2 \int_{|t|}^{\infty} f_{t}(u; \nu)\,du$.
    \item \textbf{Theoretical Rationale}: By the Central Limit Theorem, at $N = 100$ the sampling distribution of $\bar{d}$ converges closely to a Gaussian distribution $\mathcal{N}(\mu_d, \sigma_d^2 / N)$, making the paired $t$-test the uniformly most powerful unbiased test for mean makespan reduction. Against Fair FCFS + S-Shape, we obtain $t = \TStatSShape$ ($p = \PValExactSShape$, i.e., $p \PValSShape$); against Fair FCFS + Static TSP, $t = \TStatTSP$ ($p = \PValExactStaticTSP$, i.e., $p \PValStaticTSP$).
\end{itemize}

\subsubsection{3. Non-Parametric Wilcoxon Signed-Rank Test}
\begin{itemize}[leftmargin=*]
    \item \textbf{Mathematical Formulation}: Warehouse congestion delays can exhibit positive skewness or heavy tails when rare multi-picker bottlenecks occur. To ensure our conclusions do not rely on Gaussian distributional assumptions or outlier sensitivity, we perform the non-parametric Wilcoxon signed-rank test:
    \begin{enumerate}
        \item Discard exact zero differences ($d_i = 0$, ties), leaving $N_r \le N$ non-zero pairs.
        \item Sort the absolute differences $|d_i|$ in ascending order and assign ranks $R_i \in \{1, 2, \dots, N_r\}$ (averaging tied ranks).
        \item Compute the positive and negative rank sums and the test statistic $W$:
        \begin{equation}
            W^+ = \sum_{i: d_i > 0} R_i, \qquad W^- = \sum_{i: d_i < 0} R_i, \qquad W = \min(W^+, W^-).
        \end{equation}
        \item Under $H_0$ (the distribution of $d_i$ is symmetric about zero), the exact null expectation and standard deviation of $W$ are:
        \begin{equation}
            \mu_W = \frac{N_r(N_r + 1)}{4}, \qquad \sigma_W = \sqrt{\frac{N_r(N_r + 1)(2N_r + 1)}{24}}, \qquad z_W = \frac{W - \mu_W}{\sigma_W}.
        \end{equation}
    \end{enumerate}
    \item \textbf{Theoretical Rationale}: Because $W$ depends solely on the ordinal ranks $R_i$ rather than raw minute magnitudes, a single extreme wave cannot distort the test. For CAST-Wave vs.\ Fair FCFS + S-Shape ($N_r = 93$ non-tied waves, $\mu_W = 2{,}185.5$), the observed positive rank sum is just $W = \WilcoxonWSShape$ ($p = \WilcoxonPValExactSShape$, i.e., $p \WilcoxonPValSShape$), confirming overwhelming ordinal dominance. Conversely, for Variant~4 (Sector + Static TSP) vs.\ Variant~5 (Full CAST-Wave) on sparse 23-pick waves ($N_r = 31$ non-tied waves), $W = \WilcoxonWSectorVsCAST$ ($p = \WilcoxonPValSectorVsCAST > 0.05$), rigorously proving that the small average difference ($5.37\text{ min}$ vs.\ $5.32\text{ min}$) is not statistically significant and that 5-zone sectoring accounts for $\SectorShareOfGainPct\%$ of the total gain.
\end{itemize}

\subsubsection{4. Standardized Paired Effect Size (Cohen's $d_z$)}
\begin{itemize}[leftmargin=*]
    \item \textbf{Mathematical Formulation}: While $p$-values test whether an effect exists ($\mu_d \neq 0$), they scale with $\sqrt{N}$ and do not measure practical operational magnitude. We quantify standardized magnitude via Cohen's paired effect size $d_z$:
    \begin{equation}
        d_z = \frac{\bar{d}}{s_d} = \frac{t}{\sqrt{N}}.
    \end{equation}
    \item \textbf{Theoretical Rationale}: Dividing the mean reduction $\bar{d}$ by the standard deviation of the paired differences $s_d$ expresses the improvement in units of natural wave-to-wave variability. Under Cohen's canonical scale ($|d_z| = 0.20$ small, $0.50$ medium, $0.80$ large, $\ge 1.20$ very large), CAST-Wave achieves $d_z = \CohensDSShape$ vs.\ S-Shape, $d_z = \CohenDLargestGap$ vs.\ Largest-Gap, and $d_z = \CohensDTSP$ vs.\ Static TSP---confirming that the makespan reduction exceeds $1.34$ standard deviations across all three baselines.
\end{itemize}

\subsubsection{5. Empirical Dominance Analysis (Wins / Losses / Ties)}
\begin{itemize}[leftmargin=*]
    \item \textbf{Mathematical Formulation}: Using numerical tolerance $\epsilon = 10^{-6}\text{ min}$:
    \begin{equation}
        N_{\text{win}} = \sum_{i=1}^N \mathbb{I}(d_i < -\epsilon), \quad N_{\text{loss}} = \sum_{i=1}^N \mathbb{I}(d_i > \epsilon), \quad N_{\text{tie}} = \sum_{i=1}^N \mathbb{I}(|d_i| \le \epsilon).
    \end{equation}
    \item \textbf{Managerial Rationale}: Floor managers require day-to-day dispatch reliability, not merely a favorable mean driven by a subset of waves. Recording $\WinsVsSShape$ wins, $\LossesVsSShape$ losses, and $\TiesVsSShape$ ties out of $\BenchmarkWaves$ waves proves that CAST-Wave matches or outperforms the baseline in $97\%$ of historical shifts.
\end{itemize}

% ============================================================================
% SECTION 3: BLOCK-BY-BLOCK CODE EXPLANATION
% ============================================================================
\newpage
\section{Block-by-Block Codebase Walkthrough}

\subsection{Block 1: Stage 1 Multi-Dataset Raw Hygiene and Auditing}
\noindent\textbf{Source Module}: \texttt{src/clean\_data.py} (invoked via root \texttt{clean\_data.py})\\[0.5em]

\subsubsection{Operational Purpose and Handover Architecture}
As documented in \texttt{data/00\_Handover/IDEA15.5\_Dataset\_Handover.md} and \texttt{CLEANING\_LOG.md}, the raw data package in \texttt{data/01\_raw/} comprises four open-science intralogistics datasets collected from 2025--2026. Raw industrial exports exhibit mixed semicolon and comma delimiters, UTF-8 byte-order marks (BOM), Windows double-carriage-return line endings (\verb|\r\r\n|), European decimal commas, trailing whitespace in location keys, and unflagged duplicate transactions. \texttt{clean\_data.py} implements non-destructive raw-to-clean auditing across all four datasets:

\begin{enumerate}[leftmargin=*]
    \item \textbf{Dataset 1 --- Footwear Warehouse WMS (\texttt{clean\_footwear()})}:
    \begin{itemize}
        \item \textit{Exact 9-Field Deduplication}: Reads 122,370 raw order lines from \texttt{Customer\_Order.csv}, normalizes column names to \texttt{snake\_case}, and removes 2,559 exact 9-field duplicate rows, producing \TotalOrderLines~clean order lines across \TotalOrders~orders. This deduplication narrows the discrepancy between ordered units (215,460) and picked units (\TotalPicks~in \texttt{Picking\_Wave.csv}) from 13,003 units down to just \textbf{\UnitGapReconciled~units} ($97.9\%$ reconciliation).
        \item \textit{Non-Destructive Quality Flagging and Size Scale Classification}: Tags 29 zero-quantity rows and 16 missing-size rows in \texttt{quality\_flag} rather than silently deleting them, and classifies \texttt{Size (US)} into \texttt{us} (80,492 lines), \texttt{us\_x10} half-sizes (38,598 lines), \texttt{mm} (705 lines), and \texttt{unknown} (16 lines).
        \item \textit{Storage Slot Scaling and Missing Location Audit}: Parses \TotalLocations~storage slots from \texttt{Storage\_Location.csv}, scales raw coordinates by $0.08\text{ m/unit}$ into metric $(x_m, y_m)$, and cross-references \texttt{Picking\_Wave.csv} to flag 22 unlisted location codes (\texttt{quality\_flag='location\_not\_in\_slots'}).
        \item \textit{Automatic Navigation Graph Construction}: Parses 44 corridor support points from \texttt{Support\_Points\_Navigation.csv} across \texttt{LC}, \texttt{CC} (with I/O depot at \texttt{CC-08}), and \texttt{RC}, and automatically builds the 68 bidirectional corridor and cross-aisle edges in \texttt{nav\_edges.csv}.
        \item \textit{Storage Policy Melting}: Reshapes all four wide 18-column storage assignment tables (\texttt{Class\_Based}, \texttt{Dedicated}, \texttt{Random}, \texttt{Hybrid}) into tidy long-format CSVs.
    \end{itemize}

    \item \textbf{Dataset 2 --- KIT Manual Order-Picking Benchmark (\texttt{clean\_kit()})}:
    Samples 50 representative instances (5 LHS parameter sets $\times$ 10 replications) from the 13,200 shift JSON files, flattening 49,615 order-item rows into \texttt{kit\_scenarios\_flat.csv}, and normalizes decimal commas across 1,591 OFAT baseline result rows into \texttt{kit\_results\_flat.csv}.

    \item \textbf{Dataset 3 --- Fraunhofer IML DaRA Timing Calibration (\texttt{clean\_dara()})}:
    Filters the 540 raw rows of \texttt{WMS\_\_movement\_data.csv} to the 270 rack stock-out picks ($\text{Moved Quantity} < 0$) across 18 human subjects, sorts chronologically per subject, filters out session breaks ($\ge 600\text{ s}$), and computes \DaRAValidGaps~valid inter-pick gaps (median $\DaRAMedianGap\text{ s}$, mean $\DaRAMeanGap\text{ s}$, IQR $\DaRAIQR\text{ s}$). Using the literature 65\%/35\% walk-to-handling split (Chackelson et al., 2013), it exports \texttt{picker\_timing.csv} with median handling dwell $\tau_{\text{pick}} = \CalibratedPickDwell\text{ s}$ (IQR: $\DaRAHandleIQR\text{ s}$).

    \item \textbf{Dataset 4 --- Robotic Mobile Fulfillment Systems (\texttt{clean\_rmfs()})}:
    Cleans the three RMFS pod layouts (\texttt{large}: 205 nodes; \texttt{enlarged\_08} and \texttt{enlarged\_09}: 409 nodes) and melts the scientific-notation pod inventory matrices into tidy long tables (3,342, 7,016, and 7,552 non-zero pod-cargo entries). Because RMFS models Kiva-style robotic pod-to-station transport rather than manual picker-to-parts aisles, these cleaned tables are archived in \texttt{data/clean/} for completeness while the digital twin focuses on the manual footwear facility.
\end{enumerate}

% ----------------------------------------------------------------------------
\subsection{Block 2: Stage 2 Coordinate Reconciliation and Parquet Indexing}
\noindent\textbf{Source Module}: \texttt{src/data\_preprocessing.py}\\[0.5em]

\subsubsection{Resolution of the 22 Flagged Locations and Parquet Serialization}
\texttt{src/data\_preprocessing.py} ingests the Stage~1 cleaned CSVs from \texttt{data/clean/}, resolves the open modelling decisions documented in \texttt{CLEANING\_LOG.md}, and compiles typed Parquet tables into \texttt{data/02\_processed/}:
\begin{enumerate}[leftmargin=*]
    \item \textbf{Two-Rule Physical Imputation for the 22 Unlisted Pick Locations}:
    Cross-referencing the \TotalPicks~picks against the \TotalLocations~master slots identifies 22 unlisted codes: 16 staging/consolidation buffers (\texttt{RC-01}..\texttt{RC-05}, \texttt{EX-01}..\texttt{EX-11}, \texttt{C-014}..\texttt{C-102}, and \texttt{ZN-COR}) and 6 extended rack bay codes (\texttt{E-25-20}, \texttt{G-25-10}, \texttt{I-25-20}, \texttt{J-25-20}, \texttt{K-25-20}, \texttt{K-26-20}, totaling 50 picks out of \TotalPicks). Rather than dropping valid pick workload from historical waves:
    \begin{lstlisting}[language=Python]
# Rule 1: 16 Staging / consolidation codes (RC-*, ZN-*, EX-*, C-*) -> Central I/O Staging (CC-08)
is_staging = picks["x_m"].isna() & (
    picks["location"].str.startswith("RC-") | picks["location"].str.startswith("ZN-")
    | picks["location"].str.startswith("EX-") | picks["location"].str.startswith("C-")
)
picks.loc[is_staging, ["x_m", "y_m", "z", "aisle"]] = [32.24, 50.48, 1, "CC"]

# Rule 2: 6 Extended rack bay codes -> Nearest existing bay in the same aisle from slots_df
same_aisle = slots_df[slots_df["aisle"] == aisle_cand]
nearest_idx = (same_aisle["row"] - bay_cand).abs().idxmin()
x_raw_cand = float(same_aisle.loc[nearest_idx, "x_raw"])
y_raw_cand = float(same_aisle.loc[nearest_idx, "y_raw"])
    \end{lstlisting}
    This data-driven lookup guarantees $100\%$ valid coordinates across all \TotalPicks~picks without arbitrary coordinate formulas.
    \item \textbf{ABC Velocity Classification and Parquet Serialization}:
    Joins SKU ABC classes from \texttt{products\_clean.csv} onto \texttt{orders\_clean.parquet} (\TotalOrderLines~rows) and \texttt{picks\_clean.parquet} (\TotalPicks~rows), alongside \texttt{slots.parquet} (\TotalLocations~rows), \texttt{nav\_nodes.parquet} (44 rows), \texttt{nav\_edges.parquet} (68 rows), and \texttt{picker\_timing.parquet} (270 rows).
\end{enumerate}

% ----------------------------------------------------------------------------
\subsection{Block 3: Topological Navigation Graph Engine}
\noindent\textbf{Source Module}: \texttt{src/digital\_twin/warehouse\_graph.py}\\[0.5em]

\texttt{WarehouseGraph} constructs the NetworkX graph $\mathcal{G} = (\mathcal{V}, \mathcal{E})$ ($|\mathcal{V}| = 44, |\mathcal{E}| = 68$), precomputes all-pairs Dijkstra shortest paths on initialization, and exposes $O(1)$ lookup methods \texttt{distance\_to\_depot(x, y, aisle)} and \texttt{shortest\_path\_distance(p1, p2)} implementing Equations~(2)--(3).

% ----------------------------------------------------------------------------
\subsection{Block 4: Continuous-Time Discrete-Event Simulation Engine}
\noindent\textbf{Source Module}: \texttt{src/digital\_twin/simulation\_engine.py}\\[0.5em]

\texttt{WarehouseDigitalTwin} executes the event-driven state machine over \texttt{START\_BATCH}, \texttt{START\_PICK}, \texttt{FINISH\_PICK}, and \texttt{RETURN\_DEPOT} transitions, tracking active sub-aisle occupancy on $(\text{aisle}, \text{west/east})$ and computing exact wave makespan, travel distance, cumulative blocking time, and workload CV according to Equations~(4)--(8).

% ----------------------------------------------------------------------------
\subsection{Block 5: Classical and Audited Baseline Heuristics}
\noindent\textbf{Source Module}: \texttt{src/optimization/baselines.py}\\[0.5em]

Implements (i) Legacy Single-Picker Greedy FCFS (\texttt{fair\_round\_robin=False}, $P_{\text{eff}}=1, \text{CV}=\WorkloadCVLegacy$), (ii) Fair Multi-Picker Round-Robin FCFS (\texttt{batch\_orders\_fcfs\_round\_robin}, $\text{CV}=\WorkloadCVFairRR$), (iii) Balanced LPT Unit Batching (\texttt{batch\_orders\_balanced\_lpt}), and (iv) Hall's (1993) S-Shape, Largest-Gap, and Nearest-Neighbor + 2-Opt Static TSP routing.

% ----------------------------------------------------------------------------
\subsection{Block 6: Proposed CAST-Wave Spatiotemporal Optimizer}
\noindent\textbf{Source Module}: \texttt{src/optimization/cast\_wave\_optimizer.py}\\[0.5em]

Implements the three-stage architecture formalized in Section~\ref{sec:theory}: \texttt{get\_spatial\_sector()} (Equation~10), \texttt{batch\_orders()} (Equation~11), and \texttt{route\_batch\_congestion\_aware()} with shared sub-aisle reservation timeline $\mathcal{O}(a, b)$ and congestion-cost-aware 2-Opt local search (Equations~12--13, $\lambda = 15.0\text{ m}$).

% ----------------------------------------------------------------------------
\subsection{Block 7: Evaluation, Ablation, and Dynamic LaTeX Binding}
\noindent\textbf{Source Modules}: \texttt{src/evaluation/benchmark\_runner.py}, \texttt{run\_component\_ablation.py}, \texttt{generate\_visualizations.py}\\[0.5em]

\begin{enumerate}[leftmargin=*]
    \item \textbf{Unified 100-Wave Benchmark (\texttt{experiments/benchmark\_results.csv})}: Evaluates \BenchmarkWaves~historical waves ($P=4$, \texttt{seed=42}) across all policies and measures per-wave wall-clock execution time (\texttt{runtime\_ms}) via \texttt{time.perf\_counter()}.
    \item \textbf{6-Stage Component Ablation (\texttt{experiments/component\_ablation\_results.csv})}: Evaluates Variants 0 through 5 on the exact same \BenchmarkWaves~waves, logging per-stage makespan, workload CV, fleet distance, and blocked time.
    \item \textbf{3-Arm Dense-Wave Fleet Scalability (\texttt{experiments/congestion\_ablation\_results.csv})}: Evaluates $P \in \{2, 4, 6, 8, 10, 12\}$ across \DenseWavesCount~dense operational waves (averaging $\DenseMeanPicks$~picks/wave, recorded in column \texttt{num\_picks}) across \texttt{FCFS\_StaticTSP}, \texttt{Sector\_StaticTSP}, and \texttt{CAST\_Wave}.
    \item \textbf{Zero-Hardcoding Macro Generator (\texttt{generate\_visualizations.py})}: Reads the CSV artifacts directly, dynamically computes all p-value bounds and scientific strings, renders \texttt{fig1\_warehouse\_topology.pdf} and \texttt{fig2\_benchmark\_and\_scalability.pdf}, and writes \texttt{paper/tables/macros.tex} and Tables~1--3.
\end{enumerate}

% ----------------------------------------------------------------------------
\subsection{Block 8: \AutomatedTestCount-Gate Automated Verification Suite}
\noindent\textbf{Source Module}: \texttt{tests/test\_pipeline.py}\\[0.5em]

The automated \texttt{pytest} suite enforces \AutomatedTestCount~quality and regression gates:
\begin{enumerate}[leftmargin=*]
    \item \texttt{test\_clean\_data\_exists}: Verifies all cleaned CSVs and JSON audit reports exist in \texttt{data/clean/}.
    \item \texttt{test\_teammate\_cleaning\_pipeline\_outputs}: Verifies all four Stage~1 JSON cleaning reports (\texttt{footwear}, \texttt{kit}, \texttt{dara}, \texttt{rmfs}) and \texttt{CLEANING\_LOG.md} match audited row counts.
    \item \texttt{test\_processed\_data\_integrity}: Confirms exact row counts (\TotalLocations~slots, \TotalOrderLines~orders, \TotalPicks~picks, 44 nodes, 68 edges) and zero null coordinates in \texttt{data/02\_processed/}.
    \item \texttt{test\_dara\_timing\_calibration}: Verifies 270 stock-out picks, \DaRAValidGaps~valid gaps, $\DaRAMedianGap\text{ s}$ median gap, and $\CalibratedPickDwell\text{ s}$ median handling dwell.
    \item \texttt{test\_graph\_connectivity\_and\_routing}: Verifies single-component connectivity ($|V|=44, |E|=68$), zero distance at depot \texttt{CC-08}, and distance symmetry.
    \item \texttt{test\_digital\_twin\_kinematics}: Verifies positive travel times, exact pick dwell accumulation, and zero self-blocking for a single picker.
    \item \texttt{test\_cast\_wave\_optimizer\_constraints}: Confirms task conservation and strict adherence to trolley capacity $C = 27$.
    \item \texttt{test\_macro\_consistency\_with\_benchmarks}: Verifies numerical identity ($< 0.01$ rounding tolerance) between CSV outputs and \texttt{paper/tables/macros.tex}.
    \item \texttt{test\_bibliography\_and\_ieee\_template\_compliance}: Verifies 15--20 references, $\ge 80\%$ post-2021 citations, and zero undefined citation keys in \texttt{paper/main.tex}.
    \item \texttt{test\_baseline\_fair\_round\_robin\_and\_lpt\_batching}: Confirms fair round-robin and LPT load allocation across all pickers.
    \item \texttt{test\_two\_picker\_aisle\_congestion\_slowdown}: Confirms simultaneous occupancy in the same sub-aisle triggers congestion delay.
    \item \texttt{test\_zero\_credential\_leakage}: Scans all Python files for accidental secrets or API tokens.
    \item \texttt{test\_simulator\_determinism\_with\_seed}: Confirms bit-exact repeatability across runs.
    \item \texttt{test\_two\_opt\_monotonic\_improvement}: Verifies monotonic tour improvement on both baseline and CAST-Wave routers.
    \item \texttt{test\_sector\_and\_cast\_wave\_blocking\_reduction\_on\_dense\_waves}: Verifies blocking reductions across $P \in \{4, 6, 8, 10, 12\}$.
\end{enumerate}

% ============================================================================
% SECTION 4: BASELINE AUDIT AND COMPONENT ABLATION
% ============================================================================
\newpage
\section{Reviewer Audit, Baseline Dissection, and Component Ablation}

\subsection{Component-by-Component Ablation Study}
To isolate the marginal contributions of (i) route sequencing, (ii) unit workload balancing, (iii) spatial sector clustering, and (iv) dynamic congestion aversion, Table~\ref{tab:doc_component_ablation} reports the 5-stage ablation across all \BenchmarkWaves~historical waves ($P = 4$ pickers, DaRA $\CalibratedPickDwell\text{ s}$ empirical dwell).

\begin{table*}[h!]
\centering
\caption{Component-by-Component Ablation Across \BenchmarkWaves~Historical Waves ($P = 4$ Pickers, DaRA $\CalibratedPickDwell\text{ s}$ Dwell)}
\label{tab:doc_component_ablation}
\footnotesize
\resizebox{\textwidth}{!}{%
\begin{tabular}{lcccccc}
\toprule
\textbf{Ablation Variant} & \textbf{Batching Logic} & \textbf{Routing Policy} & \textbf{Makespan (min)} & \textbf{Fleet Dist. (m)} & \textbf{Blocked (s)} & \textbf{Workload CV / $\Delta$ Base} \\
\midrule
0. Legacy Single-Picker FCFS & Greedy Fill ($P_{\text{eff}}=1$) & S-Shape & $\LegacyMakespanMin \pm \LegacyMakespanStd$ & 449.9 & 0.0 & $\WorkloadCVLegacy$ (+\LegacyIncreaseVsFairPct\%) \\
1. Fair Multi-Picker FCFS (Base) & Round-Robin ($P=4$) & S-Shape & $\MakespanSShapeMin \pm \SShapeMakespanStd$ & \DistSShapeM & \BlockedSShapeSec & $\WorkloadCVFairRR$ (---) \\
2. Fair FCFS + Static TSP & Round-Robin ($P=4$) & Static 2-Opt TSP & $\MakespanTSPMin \pm \StaticTSPMakespanStd$ & \DistTSPM & \BlockedTSPSec & $\WorkloadCVFairTSP$ (-\StaticTSPRedVsSShapePct\%) \\
3. Balanced Batching + Static TSP & LPT Unit Balancing ($C=27$) & Static 2-Opt TSP & $\MakespanTSPMin \pm \StaticTSPMakespanStd$ & \DistTSPM & \BlockedTSPSec & $\WorkloadCVLPT$ (-\StaticTSPRedVsSShapePct\%) \\
4. Sector Clustering + Static TSP & 5-Zone Sector + LPT & Static 2-Opt TSP & $\MakespanSectorMin \pm \MakespanSectorStd$ & \DistSectorM & \BlockedSectorSec & $\WorkloadCVSector$ (-\SectorReductionVsFairPct\%) \\
\textbf{5. Full CAST-Wave (Proposed)} & \textbf{5-Zone Sector + LPT} & \textbf{Congestion 2-Opt ($\lambda=15$m)} & $\mathbf{\MakespanCASTMin \pm \CASTWaveMakespanStd}$ & \textbf{\DistCASTM} & \textbf{\BlockedCASTSec} & $\mathbf{\WorkloadCVCAST}$ (\textbf{-\MakespanRedSShapePct\%}) \\
\bottomrule
\end{tabular}%
}
\end{table*}

\subsection{Key Empirical Findings from the Ablation}
\begin{enumerate}[leftmargin=*]
    \item \textbf{Heuristic Sequencing (S-Shape vs.\ 2-Opt TSP)}: Replacing S-Shape traversal with 2-Opt metric graph TSP under identical round-robin batching yields a modest $\StaticTSPRedVsSShapePct\%$ makespan reduction ($\MakespanSShapeMin \rightarrow \MakespanTSPMin\text{ min}$, $\text{CV} = \WorkloadCVFairTSP$). While 2-Opt shortens within-aisle travel ($\DistSShapeM \rightarrow \DistTSPM\text{ m}$), it cannot prevent cross-warehouse travel when pickers are assigned unclustered items.
    \item \textbf{Dominance of Spatial Sector Clustering and Workload CV Trade-Off}: Incorporating 5-zone topological sector decomposition (Variant~4) slashes makespan from $\MakespanTSPMin\text{ min}$ to $\MakespanSectorMin \pm \MakespanSectorStd\text{ min}$ (a $\SectorReductionVsFairPct\%$ reduction vs.\ fair baseline, accounting for $\SectorShareOfGainPct\%$ of the total makespan reduction) and cuts fleet travel distance to $\DistSectorM\text{ m}$. Relative to unconstrained round-robin ($\text{CV} = \WorkloadCVFairRR$), sector clustering exhibits a slightly higher tour-duration dispersion ($\text{CV} = \WorkloadCVSector$) because tasks are grouped by geographic zone rather than strictly equalized across corridors, trading a small increase in duration variance for large reductions in travel distance and aisle blocking.
    \item \textbf{Role of Dynamic Congestion-Aware 2-Opt Routing}: Full CAST-Wave (Variant~5, $\lambda = 15.0\text{ m}$) achieves $\MakespanCASTMin \pm \CASTWaveMakespanStd\text{ min}$ ($-\MakespanRedSShapePct\%$ vs.\ fair baseline) and reduces narrow-aisle blocking from $\BlockedSectorSec\text{ s}$ to $\BlockedCASTSec\text{ s}$ ($-\BlockedTimeReductionPct\%$ vs.\ S-Shape). In direct pairwise comparison against Variant~4 on these moderate-size waves, CAST-Wave achieves \WinsVsSector~wins, \TiesVsSector~ties, and \LossesVsSector~losses ($W = \WilcoxonWSectorVsCAST, p = \WilcoxonPValSectorVsCAST$, not statistically significant at $\alpha = 0.05$), confirming that spatial sectoring provides the primary efficiency gain while dynamic congestion routing acts as a secondary, non-degrading operational guardrail.
\end{enumerate}

% ============================================================================
% SECTION 5: MODEL SELECTION RATIONALE
% ============================================================================
\newpage
\section{Model Selection Rationale: Why CAST-Wave and Not Others?}

Selecting the optimal algorithmic paradigm for multi-picker warehouse dispatch requires balancing four conflicting industrial requirements:
\begin{enumerate}
    \item \textbf{Sub-second Computational Latency}: Modern WMS wave dispatch operates on dynamic planning cycles, requiring order grouping and tour solutions in under 1 second.
    \item \textbf{Narrow-Aisle Congestion Awareness}: Multi-picker fleets in narrow aisles cannot pass one another without severe speed degradation.
    \item \textbf{Avoidance of Single-Picker Starvation}: A wave's makespan is determined by its \emph{straggler picker} ($\max_p T_p$).
    \item \textbf{Zero Commercial Solver Licensing Cost}: Must run natively in standard Python environments without expensive proprietary MILP licenses.
\end{enumerate}

\subsection{Formal Paradigm Comparison Matrix}

Table~\ref{tab:model_comparison} systematically compares CAST-Wave against competing optimization paradigms.

\begin{table*}[h!]
\centering
\caption{Comprehensive Comparison of Warehouse Routing and Dispatch Paradigms}
\label{tab:model_comparison}
\footnotesize
\setlength{\tabcolsep}{4pt}
\begin{tabularx}{\textwidth}{lcccccX}
\toprule
\textbf{Paradigm} & \textbf{Runtime} & \textbf{Congestion} & \textbf{Spatial} & \textbf{Stochastic} & \textbf{Optimality} & \textbf{Primary Industrial Characteristic / Failure Mode} \\
& \textbf{per Wave} & \textbf{Aversion} & \textbf{Clustering} & \textbf{Resilience} & \textbf{Guarantee} & \\
\midrule
\textbf{MILP (Exact Solvers)} & 100--10,000\,s & Explicit & Explicit & Poor & Global & Exponential time-indexed complexity explosion; cannot scale past 15 picks in real time. \\
\addlinespace[0.3em]
\textbf{Multi-Agent RL (MARL)} & 50\,ms (inf.) & Learned & Implicit & Moderate & None & Spatial transfer overfitting; policy collapses upon rack/aisle layout or ABC slotting change. \\
\addlinespace[0.3em]
\textbf{Genetic Algorithms (GA)} & 15--180\,s & Implicit & Moderate & Moderate & Near-Opt. & High runtime; requires thousands of simulation-in-the-loop fitness evaluations per wave. \\
\addlinespace[0.3em]
\textbf{Classical S-Shape / Gap} & \SShapeRuntimeMs\,ms & \textbf{None} & \textbf{None} & High & Heuristic & Single-picker assumption; causes $\BlockedSShapeSec\text{ s}$ aisle blocking and $\DistSShapeM\text{ m}$ travel per wave. \\
\addlinespace[0.3em]
\textbf{Static 2-Opt TSP} & \StaticTSPRuntimeMs\,ms & \textbf{None} & \textbf{None} & High & Heuristic & Optimizes isolated tour length ($\DistTSPM\text{ m}$) but leaves $\BlockedTSPSec\text{ s}$ of multi-picker aisle blocking. \\
\addlinespace[0.3em]
\rowcolor{secondary!10}
\textbf{CAST-Wave (Proposed)} & \textbf{\CASTWaveRuntimeMs\,ms} & \textbf{Explicit} & \textbf{Explicit} & \textbf{High} & \textbf{High} & \textbf{Achieves $-\MakespanRedSShapePct\%$ makespan, $-\DistanceImprovementVsSShape\%$ distance, and $-\BlockedTimeReductionPct\%$ blocking vs.\ fair baseline.} \\
\bottomrule
\end{tabularx}
\end{table*}

\subsection{Deep-Dive Analysis of Alternative Paradigms}

\subsubsection{1. Why Not Exact Mixed-Integer Linear Programming (MILP / Branch-and-Cut)?}
\begin{rationalebox}{The MILP Scalability Barrier}
An exact mathematical formulation of multi-picker routing with continuous-time collision avoidance requires time-indexed binary decision variables $x_{i, j, p, t} \in \{0, 1\}$ subject to mutual exclusion sub-aisle constraints $\sum_{p=1}^P \sum_{\tau = t}^{t + \Delta t} x_{u, v, p, \tau} \le K_{\text{aisle}} = 1$.
\end{rationalebox}
\textbf{Failure Rationale}: The decision space scales as $O(|V|^2 \cdot P \cdot T_{\max})$, creating millions of binary variables for a 50-item wave and requiring hundreds of seconds in commercial branch-and-cut solvers (Briant et al., 2025), failing real-time WMS dispatch deadlines.

\subsubsection{2. Why Not Multi-Agent Deep Reinforcement Learning (MARL)?}
\begin{rationalebox}{The MARL Spatial Generalization Flaw}
Multi-Agent Reinforcement Learning models (e.g., MAPPO, QMIX) train neural network agents to coordinate movements via shared value functions or decentralized actor policies.
\end{rationalebox}
\textbf{Failure Rationale}: Deep RL policies overfit to specific coordinate topologies and require retraining whenever seasonal footwear SKUs are re-slotted across ABC zones, while offering no deterministic worst-case tour bounds.

\subsubsection{3. Why Not Pure Metaheuristics (Genetic Algorithms / Simulated Annealing)?}
\begin{rationalebox}{The Metaheuristic Simulation-in-the-Loop Bottleneck}
Genetic Algorithms optimize chromosome representations of pick sequences by evaluating population fitness across successive generations.
\end{rationalebox}
\textbf{Failure Rationale}: Evaluating narrow-aisle congestion inside a GA fitness function requires thousands of discrete-event simulations per wave ($N_{\text{pop}} \times N_{\text{gen}}$), inflating runtime from milliseconds to minutes.

\subsubsection{4. Why Not Unclustered Static 2-Opt TSP?}
\begin{rationalebox}{The Unclustered Routing Bottleneck}
Static 2-Opt TSP minimizes walking distance $D = \sum d(i, i+1)$ for each picker's assigned batch in isolation.
\end{rationalebox}
\textbf{Failure Rationale}: Without spatial sector clustering during batch formation, each picker's round-robin batch spans all 18 aisles of the warehouse. Even with 2-Opt refinement, the fleet walks $\DistTSPM\text{ m}$ and suffers $\BlockedTSPSec\text{ s}$ of aisle blocking ($\MakespanTSPMin\text{ min}$ makespan), whereas CAST-Wave's sector-cohesive batching and congestion-aware 2-Opt cut fleet distance to $\DistCASTM\text{ m}$ and blocking to $\BlockedCASTSec\text{ s}$ ($\MakespanCASTMin\text{ min}$ makespan).

% ============================================================================
% SECTION 6: MATHEMATICAL AND EXPERIMENTAL REFERENCE
% ============================================================================
\section{Algorithmic Parameters and Validation Summary}

\subsection{Calibrated Physical and Simulation Constants}
Table~\ref{tab:constants} summarizes the empirical constants calibrated into the digital twin.

\begin{table}[h!]
\centering
\caption{Calibrated Industrial Digital Twin Parameters}
\label{tab:constants}
\resizebox{\textwidth}{!}{%
\begin{tabular}{llcl}
\toprule
\textbf{Parameter} & \textbf{Symbol} & \textbf{Value} & \textbf{Empirical Provenance / Basis} \\
\midrule
Nominal Walk Speed & $v_0$ & $\CalibratedWalkSpeed\text{ m/s}$ & Industrial picker velocity standard (Franzke et al., 2017) \\
Pick Dwell Handling Time & $\tau_{\text{pick}}$ & $\CalibratedPickDwell\text{ s}$ & Fraunhofer DaRA 2026 empirical median handling time (IQR: \DaRAHandleIQR~s) \\
Sub-Aisle Unimpeded Capacity & $K_{\text{aisle}}$ & $1\text{ picker}$ & Narrow sub-aisle $(\text{aisle}, \text{west/east})$ single-cart clearance \\
Congestion Slowdown Factor & $\alpha$ & $0.40$ & Velocity degradation under simultaneous sub-aisle occupancy \\
Trolley Unit Capacity & $C$ & $27\text{ units}$ & Physical footwear picking trolley capacity limit \\
Congestion Penalty Coeff. & $\lambda$ & $15.0\text{ m}$ & Calibrated sub-aisle collision penalty ($\approx 13.6\text{ s}$ yielding delay) \\
Depot Node Identifier & Depot & \texttt{CC-08} & Central corridor I/O point at $(x=32.24\text{ m}, y=50.48\text{ m})$ \\
\bottomrule
\end{tabular}%
}
\end{table}

\subsection{Computational Complexity and Measured CPU Latency}
Table~\ref{tab:complexity} provides the formal asymptotic complexity and empirically measured wall-clock latency (\texttt{time.perf\_counter()} across \BenchmarkWaves~waves) for all functional modules.

\begin{table}[h!]
\centering
\caption{Asymptotic Complexity and Empirically Measured CPU Latency}
\label{tab:complexity}
\resizebox{\textwidth}{!}{%
\begin{tabular}{llll}
\toprule
\textbf{Pipeline Module} & \textbf{Time Complexity} & \textbf{Space Complexity} & \textbf{Measured CPU Latency} \\
\midrule
Stage 1 Multi-Dataset Cleaning (\texttt{clean\_data.py}) & $O(N_{\text{raw}})$ & $O(N_{\text{raw}})$ & $\approx 9.0\text{ s}$ (offline, 4 datasets) \\
Stage 2 Coordinate \& Parquet Prep (\texttt{data\_preprocessing.py}) & $O(N \log N)$ & $O(N)$ & $\approx 1.2\text{ s}$ (offline) \\
Graph Dijkstra All-Pairs (\texttt{warehouse\_graph.py}) & $O(|V| \cdot |E| + |V|^2 \log |V|)$ & $O(|V|^2)$ & $\approx 8.0\text{ ms}$ (one-time init) \\
Fair FCFS + S-Shape Routing & $O(M \log M)$ & $O(M)$ & $\SShapeRuntimeMs\text{ ms}$ mean per wave \\
Fair FCFS + Static 2-Opt TSP & $O(P \cdot (M/P)^2)$ & $O(M)$ & $\StaticTSPRuntimeMs\text{ ms}$ mean per wave \\
CAST-Wave Sector Batching \& Congestion 2-Opt & $O(M \log M + P \cdot (M/P)^2)$ & $O(M + P)$ & $\CASTWaveRuntimeMs\text{ ms}$ mean ($\CASTWaveRuntimeMaxMs\text{ ms}$ max) \\
Discrete-Event Simulation Engine (\texttt{simulate\_wave}) & $O(M \log P)$ & $O(P + M)$ & $\approx 0.4\text{ ms}$ per wave \\
\bottomrule
\end{tabular}%
}
\end{table}

% ============================================================================
% SECTION 7: POINT-BY-POINT RESPONSE TO CODE REVIEW
% ============================================================================
\section{Technical Dissection and Response to Peer Audits}
\label{sec:review_response}

This section documents the engineering analysis and verified fixes implemented across the two independent audit reviews.

\subsection{Resolution of Router Confounding and Congestion-Aware 2-Opt}
\begin{enumerate}[leftmargin=*]
    \item Aligned the optimizer's reservation table $\mathcal{O}(a, b)$ with the simulator's exact physical sub-segments $(\text{aisle}, \text{west/east})$.
    \item Calibrated the collision penalty to $\lambda = 15.0\text{ m}$, matching the physical $12\text{--}15\text{ s}$ speed-yielding delay at $v_0 = \CalibratedWalkSpeed\text{ m/s}$ and $\alpha = 0.40$.
    \item Embedded congestion-cost-aware 2-Opt local search directly inside \texttt{route\_batch\_congestion\_aware}: edge swaps evaluate total penalized cost $J_{\text{total}}(\pi_p)$, eliminating geometric loops while preserving valid congestion detours.
\end{enumerate}

\subsection{Honest Multi-Picker Baselines and Dense-Wave Boundary Conditions}
\begin{enumerate}[leftmargin=*]
    \item \textbf{Fair Multi-Picker Round-Robin Baseline}: All primary claims benchmark against Fair Multi-Picker Round-Robin FCFS ($\MakespanSShapeMin \pm \SShapeMakespanStd\text{ min}, \text{CV} = \WorkloadCVFairRR$), achieving a verified $-\MakespanRedSShapePct\%$ makespan reduction ($t = \TStatVsSShape, p = \PValExactSShape, W = \WilcoxonWSShape, d_z = \CohensDSShape$; \WinsVsSShape~wins, \LossesVsSShape~losses, \TiesVsSShape~ties), a $-\DistanceImprovementVsSShape\%$ distance reduction, and a $-\BlockedTimeReductionPct\%$ blocking reduction.
    \item \textbf{Decoupling Sector Clustering from Dynamic Congestion Routing}: Component ablation (Table~\ref{tab:doc_component_ablation}) demonstrates that 5-zone spatial sectoring delivers $\SectorShareOfGainPct\%$ of the makespan reduction on moderate 23.4-pick waves, while the additional makespan difference of Full CAST-Wave ($5.37 \to 5.32\text{ min}$, $p = \WilcoxonPValSectorVsCAST$) is non-significant at $\alpha = 0.05$, acting as a secondary guardrail that trims mean blocking from $\BlockedSectorSec\text{ s}$ to $\BlockedCASTSec\text{ s}$.
    \item \textbf{Dense-Wave Fleet Scalability ($P = 2$ to $12$)}: Across \DenseWavesCount~merged dense waves ($\approx \DenseMeanPicks$~picks/wave), CAST-Wave achieves the lowest blocking at $P \in \{2, 4, 6, 8\}$ ($\DenseBlockedCASTTwo\text{ s}, \DenseBlockedCASTFour\text{ s}, \DenseBlockedCASTSix\text{ s}, \DenseBlockedCASTEight\text{ s}$) and lowest makespan at $P = 2$ ($\DenseMkCASTTwo\text{ min}$), while static sectoring achieves slightly lower makespan for $P \ge 4$ ($\DenseMkSectorFour\text{ min}$ vs.\ $\DenseMkCASTFour\text{ min}$ at $P=4$; $\DenseMkSectorTwelve\text{ min}$ vs.\ $\DenseMkCASTTwelve\text{ min}$ at $P=12$) when multiple pickers must share each of the five sectors. Both sector-based methods slash uncoordinated FCFS's $\DenseBlockedFCFSTwelve\text{ s}$ blocking queue by over $53\%$.
\end{enumerate}

% ============================================================================
% SECTION 8: CONCLUSION
% ============================================================================
\section{Conclusion}
The complete software and data architecture in \texttt{WarehouseDatasets2} integrates rigorous multi-dataset hygiene (\texttt{src/clean\_data.py} and \texttt{src/data\_preprocessing.py}), exact 44-node topological graph modeling, DaRA-calibrated continuous-time physics simulation, and the \textbf{CAST-Wave} optimization engine. Evaluated against fair multi-picker baselines across \BenchmarkWaves~historical waves, CAST-Wave delivers a statistically robust $-\MakespanRedSShapePct\%$ makespan reduction ($p = \PValExactSShape, d_z = \CohensDSShape$), a $-\DistanceImprovementVsSShape\%$ reduction in fleet travel distance, and a $-\BlockedTimeReductionPct\%$ reduction in narrow-aisle blocking in an empirically measured average of $\CASTWaveRuntimeMs\text{ ms}$ per wave.

\end{document}
